\section{Cyclic simplex formulas for the trace}\label{sec:cyclic}

We now express a conjugacy-class trace coefficient as a weighted average
of curvature Pfaffians. Idempotency provides compatibility of the
weights under deletion of a vertex; cyclicity supplies the cancellation
that makes the average independent of the connection form.
The tuple weights have the same algebraic origin as the weighted
paths in \cite[Section~7]{OpenAIBass2026}; we give the identities
needed here directly.

\subsection{Tuples and their weights}
Fix $g\in G$ and write $H=C_G(g)$ for its centralizer. For $k\ge0$ let
\[
 \mathcal T_k(g)=H\backslash G^{k+1},
\]
where $H$ acts by simultaneous left multiplication. For a tuple
$h=(h_0,\ldots,h_k)$ extend its indices to all integers by
$h_{i+k+1}=gh_i$ and define
\begin{equation}\label{cy:weights}
 v_{k+1}(h)=\tr\left(\prod_{i=0}^k p(h_i^{-1}h_{i+1})\right),
 \qquad l_i=1+d(h_{i-1},h_i).
\end{equation}
The product is in the displayed order. Both expressions descend to
$\mathcal T_k(g)$, and $l_i$ is periodic with period $k+1$.

\begin{lemma}[Weight identities]\label{cy:weights-lemma}
The following properties hold.
\begin{enumerate}[label=\textup{(\roman*)}]
\item The weights have total sum $\tau_{[g]}(p)$, in every dimension.
\item Twisted rotation
$(h_0,\ldots,h_k)\mapsto(h_1,\ldots,h_k,gh_0)$ preserves the weights.
\item For $k\ge1$, summing the weights over all insertions of one entry
into a fixed shortened tuple gives its weight in $\mathcal T_{k-1}(g)$.
\end{enumerate}
All these sums converge absolutely, even after multiplication by any
fixed polynomial in the edge lengths. In particular, with
$M_p=\sum_{w\in G}|p(w)|_1(1+|w|)$,
\begin{equation}\label{cy:moment}
 \sum_{h\in\mathcal T_k(g)}|v_{k+1}(h)|\prod_{i=0}^k l_i
 \le M_p^{k+1}.
\end{equation}
\end{lemma}
\begin{proof}
Put $w_i=h_i^{-1}h_{i+1}$. Their product is $h_0^{-1}gh_0$.
Conversely, any list $(w_0,\ldots,w_k)$ whose product lies in $[g]$
determines a tuple: choose $h_0$ realizing this conjugate and successively
put $h_{i+1}=h_iw_i$. Two choices of $h_0$ differ by left multiplication
by an element of $H$, and this is the only ambiguity. We therefore have
a bijection between $\mathcal T_k(g)$ and such increment lists.

The inequality
$|\tr(A_0\cdots A_k)|\le\prod_i|A_i|_1$ gives
\eqref{cy:moment} upon dropping the restriction on the product of the
increments. The exponential moment gives the same conclusion for any
fixed polynomial in their lengths. The unrestricted convolution sum
with product in $[g]$ is the trace coefficient of $p^{k+1}=p$, proving
\textup{(i)}. Rotation is cyclic permutation of the matrix factors,
so \textup{(ii)} follows from their trace.

To delete an interior entry, sum the two adjacent factors over that
entry. Their convolution is $p^2=p$. The other deletions follow by
rotation. One may fix a representative of the shortened tuple because
its simultaneous left $H$-action has trivial stabilizer. Thus no
multiplicity enters the insertion sum, and \textup{(iii)} follows.
\end{proof}

\subsection{Compatible forms}
For each tuple use the ordered time domain
\[
 \widetilde\Delta_k
 =\{(t_0,\ldots,t_k):t_0\le t_1\le\cdots\le t_k\le t_0+1\}.
\]
Extend times by $t_{i+k+1}=t_i+1$, and set
$x_i=t_{i+1}-t_i$ for $0\le i\le k$. Thus $x_i\ge0$ and
$\sum_i x_i=1$. Simultaneous time translation is generated by
$V=\sum_{i=0}^k\partial_{t_i}$. Its quotient is the simplex
\[
 \Delta_k=\{0\le t_1\le\cdots\le t_k\le1\},
\]
using the section $t_0=0$, oriented by $dt_1\wedge\cdots\wedge dt_k$.
We regard $[t_i,t_{i+1})$ as carrying the label $h_i$.
The face $x_i=0$ corresponds to deleting $(h_i,t_i)$: the interval
carrying the label $h_i$ has length zero. At the closing face use the
periodic lifts to make this same convention.

\begin{definition}\label{cy:connection}
A compatible connection family through dimension $n$ assigns a real
one-form $\alpha_h$ to each $h\in\mathcal T_k(g)$, $0\le k\le n$,
on $\widetilde\Delta_k$, with the following properties:
\begin{enumerate}[label=\textup{(\roman*)}]
\item It is invariant under time translation, and $\alpha_h(V)=1$.
\item It is invariant under simultaneous left multiplication by $H$
and under twisted rotation of both labels and times.
\item Its tangential restriction to $x_i=0$ agrees with the form of
the shortened tuple obtained by deleting $(h_i,t_i)$.
\item The form is smooth up to every face, by restriction of a smooth
extension, and the coefficients of $\alpha_h$ and $d\alpha_h$ are bounded
by polynomials in the edge lengths, for each fixed dimension and fixed
construction parameters.
\end{enumerate}
\end{definition}

Translation invariance and $\alpha(V)=1$ imply
$\iota_Vd\alpha=0$. Thus $d\alpha$ is a form on $\Delta_k$ independent
of the choice of translation section. A translation-invariant form
whose contraction with $V$ is zero will be called \emph{basic}.

For a real skew-symmetric $2m\times2m$ matrix $A$, our convention is
\[
 \omega=\sum_{i<j}A_{ij}\,dt_i\wedge dt_j
 \quad\Longrightarrow\quad
 \omega^m=m!\Pf(A)\,dt_1\wedge\cdots\wedge dt_{2m}.
\]
In particular $\Pf(A)^2=\det A$. This identity follows by reducing a
skew form to two-dimensional blocks, with both sides transforming by
the square of the determinant of a change of basis.

\begin{proposition}[Trace as a curvature average]\label{cy:trace-identity}
Let $n=2m>0$ and let $\alpha$ be a compatible connection family through
dimension $n$. Write $A_h$ for the matrix of $d\alpha_h$ on the section
$t_0=0$. If $\mathbb E_{\Delta_n}$ denotes uniform probability on
$\Delta_n$, then
\begin{equation}\label{cy:trace-formula}
 \sum_{h\in\mathcal T_n(g)}v_{n+1}(h)\,
       \mathbb E_{\Delta_n}\Pf(A_h)
   =(-1)^m\tau_{[g]}(p).
\end{equation}
The series converges absolutely.
\end{proposition}
\begin{proof}
First compare two compatible families $\alpha$ and $\alpha'$.
Their difference is basic. The following transgression form is the
elementary abelian instance of the connection-invariance mechanism in
Chern--Weil theory \cite[Sections~2--3]{ChernSimons1974}:
\[
 \eta=(\alpha-\alpha')\wedge
       \sum_{j=0}^{m-1}(d\alpha)^j\wedge(d\alpha')^{m-1-j}
\]
is basic, has degree $n-1$, and satisfies
$d\eta=(d\alpha)^m-(d\alpha')^m$. It respects all the face and
rotation identifications. Stokes' theorem and
Lemma~\ref{cy:weights-lemma}\textup{(iii)} express the weighted sum of
$\int_{\Delta_n}d\eta$ as a signed sum of copies of
\[
 I=\sum_{h\in\mathcal T_{n-1}(g)}v_n(h)
                       \int_{\Delta_{n-1}}\eta_h.
\]
Each face identification has a fixed orientation sign, independent
of the labels. The polynomial bounds and the exponential moment
justify Stokes term by term and all the rearrangements.

On the simplex of dimension $n-1$, rotation cyclically permutes its
$n$ gap coordinates. It has orientation sign $(-1)^{n-1}=-1$.
It preserves the weights and the forms. Basicness ensures that
returning to the section $t_0=0$ has no additional effect on the
integral. Consequently $I=-I$, so $I=0$. The weighted integral of
$(d\alpha)^m$ is therefore independent of the compatible family.

Use now the particular connection
\[
 \alpha^{\rm can}=\sum_{i=0}^k x_i\,dt_i.
\]
Its contraction with $V$ is one, and all other conditions are immediate
from the gap coordinates: on a zero gap the corresponding summand
vanishes and adjacent time differentials agree tangentially.
On $\widetilde\Delta_k$ its curvature is
\[
 d\alpha^{\rm can}=\sum_{i=0}^k dt_{i+1}\wedge dt_i.
\]
For $k=n$ in the section $t_0=0$, the terms at either end vanish,
leaving the path of adjacent pairs. Its Pfaffian is $(-1)^m$:
the only complete pairing uses $(1,2),(3,4),\ldots,(n-1,n)$,
each with coefficient $-1$.

Finally, $\Delta_n$ has Euclidean coordinate volume $1/n!$, so
\begin{equation}\label{cy:factorials}
 \frac{n!}{m!}\int_{\Delta_n}(d\alpha_h)^m
 =\mathbb E_{\Delta_n}\Pf(A_h).
\end{equation}
Evaluation at $\alpha^{\rm can}$ and the total-weight identity give
\eqref{cy:trace-formula}.
\end{proof}

The proposition reduces trace vanishing to a quantitative construction.
We shall choose connections, in arbitrarily large even dimensions,
for which the average of the absolute Pfaffian is at most a small
exponential factor times $\prod_{i=1}^n l_i$.
Equation~\eqref{cy:moment} then permits summation against the tuple
weights. The next section constructs the spatial summaries used to
obtain these connections.
